\section{Conclusion}\label{sec:conclusion}

How an agent behaves when its tools break is a property that can be stated, injected for, and checked. We proposed a failure model with faults on one axis and unsafe behaviours on the other, and a method that adapts chaos engineering to agents: a transparent proxy injects faults on the MCP tool path at deterministic points, deterministic predicates over the call trace and the final answer grade the result with no model judge, and repeated runs with a pass threshold and paired no-fault controls turn a stochastic agent's behaviour into a verdict. AgentChaos implements the method on switchyard, an open proxy core meant to be shared with other runtime tools. Prior work has already shown that tool faults degrade agents, has built injection middleware for MCP, and has graded deterministically. Our contribution is narrower: declared behavioural predicates over the trace and the answer, checked on real tools behind a transparent proxy; pass thresholds with paired controls, so that a result can gate a change; and a reusable proxy core.
The pilot was modest and so are its findings. On a simple single-domain task, both models handled all four faults repeatably and, on our reading of the transcripts (to be validated in RQ3), safely. The one failing verdict was a regex false positive. The pilot therefore says more about the grader than about agents. Deterministic checks of call behaviour were stable and uncontroversial, while regex checks of free-text answers were the fragile part.
The open questions are empirical: whether verdicts are stable enough, and deterministic checks faithful enough, to gate a merge. We have set out a protocol to answer them, and we invite others to use the tool, criticise the failure model, and contribute experiments.
